\section{Structural Consequences for Exact Certification}\label{sec:structural-consequences}

The regime hierarchy yields both theoretical barriers and practical implications. This section collects structural consequences that follow from the complexity classifications: witness-checking lower bounds, approximation gaps, configuration-simplification limits, regime-dependent certification competence, and the simplicity tax principle. The common theme is that exact relevance certification imposes constraints that cannot be circumvented without either restricting the regime or accepting tradeoffs.

\subsection{Exact Simplification and Over-Specification}

Many practical simplification questions are instances of sufficiency checking. For example, configuration parameter reduction asks whether a subset of parameters preserves all decision-relevant behavior, which directly reduces to \textsc{Sufficiency-Check} (Appendix~\ref{app:applications} contains detailed reductions). This connection yields immediate complexity consequences.

\leanmetapending{\LHrng{DQ}{36}{37}}
\begin{corollary}[No General-Purpose Exact Configuration Minimizer]
\label{cor:no-general-minimizer}
Assuming $\Pclass \neq \coNP$, there is no polynomial-time general-purpose procedure that takes an arbitrary succinctly encoded configuration problem and always returns a smallest behavior-preserving parameter subset.
\end{corollary}

\begin{proof}
Such a procedure would solve \textsc{Minimum-Sufficient-Set} in polynomial time for arbitrary succinctly encoded instances, contradicting Theorem~\ref{thm:minsuff-conp} under the assumption $\Pclass \neq \coNP$.
\end{proof}

\leanmetapending{\LH{CR1}, \LH{CT12}, \LH{DQ26}, \LH{HD14}}
\begin{remark}
The corollary does not rule out useful simplification tools; it rules out a fully general exact minimizer with worst-case polynomial guarantees. The structured regimes isolated in Section~\ref{sec:tractable} remain viable precisely because they restrict the source of hardness. Moreover, when certification is computationally expensive but parameter maintenance is cheap, over-specification can be cost-optimal—another manifestation of the simplicity tax principle.
\end{remark}

\subsection{Regime-Limited Exact Certification}

\leanmetapending{\LH{DQ36}, \LH{DQ40}, \LH{DQ44}}
\begin{proposition}[Exact Certification Competence Depends on Regime]
\label{prop:competence-by-regime}
Within the model of this paper, exact certification competence is regime-dependent. In the general succinct regime, exact relevance certification and exact minimization inherit the hardness results of Sections~\ref{sec:hardness} and~\ref{sec:dichotomy}. In the structured regimes of Section~\ref{sec:tractable}, exact certification becomes available in polynomial time.
\end{proposition}


\begin{proof}
The negative side is given by Theorems~\ref{thm:sufficiency-conp} and~\ref{thm:minsuff-conp}, together with the ETH-conditioned lower bounds summarized in Section~\ref{sec:dichotomy}. The positive side is given by the tractability results of Section~\ref{sec:tractable}, which provide explicit polynomial-time certification procedures under structural restrictions.
\end{proof}

The next three subsections sharpen that basic regime distinction in different directions: exact reliability under polynomial budgets, witness-checking and approximation limits, and the cost of compressing a central interface while leaving exact relevance unresolved.
